\section{Comprobación de perfiles, evaluación y eventos}\label{sec:pruebas-perfiles-plazos}

La comprobación del catálogo separa representación, reproducción matemática, autorización, persistencia y transporte. Los casos utilizan perfiles sintéticos y fuentes controladas para detectar pérdidas de precisión y cambios de referencia. Sus resultados no acreditan un perfil jurídico aplicable a cualquier expediente. El registro de ejecución del catálogo PostgreSQL y de HTTP se conserva en \rutarepo{docs/verification-report.md}; las cifras globales de las secciones precedentes conservan sus cortes históricos y no se reutilizan como resultados de estos módulos.

\subsection{Reversibilidad y resultados reproducibles}

Las pruebas de \texttt{DINP1} recorren selección conocida y desconocida, calificación opcional, familias exactas, requisitos, reglas y calendario. Verifican que una solicitud representable conserve sus valores aun cuando su evaluación posterior deba bloquearse. Los negativos alteran etiquetas, longitudes, revisiones, cantidades y componentes de tiempo, así como truncamiento y bytes sobrantes. La conservación de un UUID cero se comprueba separadamente de la ausencia de referencia.

Las pruebas de \texttt{DPRF1} contrastan ámbitos, cantidades fijas y ordenadas, políticas de corte y corpus con calendarios completos. Incluyen los límites de conteo antes de reservar memoria y la lectura estricta del calendario anidado. Un vector mínimo independiente y las alteraciones de sus campos evitan que el emisor y el lector acepten conjuntamente un mismo error de formato. Los instantes esperados conservan nanosegundos y desfase original; se ejercitan representaciones válidas del tipo de instante que no pertenecen al perfil más restringido del tiempo declarado de un acto.

La construcción del perfil reproduce el resultado esperado de cada ejemplo. Las pruebas rechazan referencias ajenas, identificadores repetidos, ejemplos incompatibles y corpus sin una candidata. Los ensayos distinguen cantidad ordenada ausente, inesperada y superior al máximo, sin rellenar o recortar el valor. También separan candidata civil, instante y bloqueos del motor. El límite superior de tamaño del perfil se usa como cota conservadora de admisión; no se presenta como un máximo alcanzable por una definición semánticamente válida.

El evaluador se comprueba con fuentes y recibos coherentes antes de introducir alteraciones. Se exige que los errores de integridad precedan a los bloqueos de aplicabilidad. Los casos conservan cálculo parcial ante condiciones pendientes y rechazan condiciones duplicadas o ajenas. Una regla no instanciada permanece ausente. Una candidata civil sin corte, fuera de su cobertura o acompañada de otro bloqueo no produce instante final. El corte usa su propio desfase y la extracción conserva la precisión del acto.

\subsection{Autorización, atomicidad y restauración}

Los ensayos del servicio diferencian autenticación, preparación y confirmación. Comprueban permisos de los cuatro roles, reautenticación, cambios de identidad y respuesta del puerto. Las consultas contrastan ámbito global y privado, orden, límites de página, cursor y revisión exacta. Una colección global no puede revelar un perfil privado mediante listado, detalle o historia; autorizar un expediente es una condición previa a su colección.

Las pruebas PostgreSQL utilizan bases desechables y el rol de ejecución. Recorren publicación, reemplazo y retiro, así como transacciones concurrentes sobre una misma base. Los rechazos por revisión obsoleta, cambio de actor, cierre o fallo de auditoría deben conservar el estado anterior. El retiro es terminal y copia la definición y el algoritmo de la revisión precedente. El autor histórico conserva su correo capturado aunque la cuenta cambie después; la ausencia de esa cuenta, en cambio, es una inconsistencia del inventario.

Los casos de catálogo dañan deliberadamente funciones, restricciones, proyecciones, disparadores y permisos para comprobar el rechazo de arranque. Incluyen una raíz con UUID cero y una definición cuya cabecera SQL es válida pero cuyo cuerpo completo es inválido. Ese caso discrimina la función de la proyección ligera frente al lector Rust. Otro caso introduce un hijo heredado vacío: que el contenido visible coincida no basta para aceptar una tabla cuya consulta podría incorporar filas externas al esquema esperado. La comprobación exige tablas permanentes y sin relaciones de herencia.

La restauración se ensaya con perfiles globales y privados, revisiones publicadas, reemplazadas y retiradas, y un autor que posteriormente quedó inactivo. El procedimiento conserva instantáneas de tablas y estado de la secuencia, realiza \texttt{pg\_dump} y \texttt{pg\_restore}, y vuelve a consultar cada revisión exacta. Una corrección posterior con otro Owner comprueba la continuidad de historia y emisión de eventos. La importación heredada debe rechazar destinos ocupados; no fusiona registros de procedencias distintas.

\subsection{Emisión de cambios y frontera HTTP}

Las pruebas de eventos verifican identidad, revisión, operación y ámbito. Distinguen expresamente el expediente ausente de un perfil global y el expediente concreto de un perfil privado. Los casos de concurrencia y fallo controlado comprueban que revisión, auditoría y evento compartan resultado transaccional, mientras la secuencia puede conservar huecos de operaciones revertidas. También comprueban la ausencia de eventos históricos inventados durante una migración y la conservación de los existentes tras restaurar.

La comprobación de transporte usa un flujo de trabajo simulado para aislar rutas, autenticación, cuerpos y proyecciones. Se rechazan campos desconocidos o repetidos, formas temporales incompatibles, revisiones inválidas y discrepancias entre ruta y comando. Los resultados esperados y calendarios completos se proyectan sin perder componentes; las páginas no exponen el corpus completo. El presupuesto del cuerpo se fundamenta en una cota estructural conservadora de JSON compacto, diferenciada de una medición de rendimiento o de una petición máxima observada.

Las pruebas focales y los recorridos con servicios reales son evidencias distintas. El contrato HTTP implementado no demuestra por sí solo la persistencia de evaluaciones, consumo de eventos, entrega de correo o usabilidad de la futura interfaz. El \cref{sec:matriz-perfiles-plazos} relaciona los archivos y fronteras comprobadas; el informe técnico conserva los resultados de ejecución correspondientes. La aceptación de reglas jurídicas y del seguimiento operativo permanece separada de esta comprobación técnica.
